\fullwidth{You may use a calculator to solve the following problem but you must \textbf{show your work} in the space provided.  You may \textbf{round to 3 decimal places after each step \emph{or} after the final calculation}. Any other rounding will result in a 3 point deduction.}

\question[15]
Whether your left or right thumb is on top when you clasp your hands together is a simple genetic trait with alleles \textit{T} (dominant) and \textit{t} (recessive).  In a previous class, 24 out of 84 students placed their right thumb on top, meaning they were homozygous for the recessive allele.  Assuming that the class was in Hardy-Weinberg equilibrium for this character, calculate the following values. The final values may sum to 0.999, 1 or 1.001 due to rounding error. (3 pts each, 15 pts total)

\vspace{2\baselineskip}
Frequency of \textit{T}:

\vspace{2\baselineskip}
Frequency of \textit{t}:

\vspace{2\baselineskip}
Frequency of \textit{TT}:

\vspace{2\baselineskip}
Frequency of \textit{Tt}:

\vspace{2\baselineskip}
Frequency of \textit{tt}: 

\begin{solution}
	\begin{multicols}{2}
\vspace{2\baselineskip}
Frequency of \textit{T}: 0.465

\vspace{2\baselineskip}
Frequency of \textit{t}: 0.535

\vspace{2\baselineskip}
Frequency of \textit{TT}: 0.216

\vspace{2\baselineskip}
Frequency of \textit{Tt}: 0.498

\vspace{2\baselineskip}
Frequency of \textit{tt}: 0.286

\vspace{\baselineskip}
Sums to 1.

	\columnbreak

\vspace{2\baselineskip}
If \textit{t} = 0.286

\vspace{2\baselineskip}
\textit{T} = 0.714

\vspace{2\baselineskip}
\textit{TT} = 0.510

\vspace{2\baselineskip}
\textit{Tt} = 0.408

\vspace{2\baselineskip}
\textit{tt} = 0.082

\vspace{\baselineskip}
Sums to 1.

\end{multicols}

\end{solution}